% ===========================================================================
\chapter{Terms, atoms and rules}\label{ch:terms}

Before the engine can read a program, we need something to read it
\emph{into}: Nex classes for the things a Datalog program is made of. They
are few.

\begin{description}[style=nextline, leftmargin=1.2em, itemsep=0.4em]
  \item[Terms.] The smallest pieces: a constant such as \code{alice}, \code{42}
    or \code{"Ann Lee"}, or a variable such as \code{X}.
  \item[Atoms.] A predicate applied to terms: \code{parent(alice, bob)},
    \code{ancestor(X, Y)}. An atom with no variables is \emph{ground}, and a
    fact is a ground atom.
  \item[Literals.] The conditions in the body of a rule. Most are atoms,
    which must match a fact. Two other kinds test something instead:
    \code{X != Y} compares two terms, and \code{not banned(U)} asks that a
    fact \emph{not} hold.
  \item[Rules.] A head and a body: the head holds whenever the whole body
    does.
\end{description}

\section{Terms}

A term is exactly one of four things, each carrying different data, which
is what a \code{union} is for. Three small functions answer the questions
the rest of the engine asks about a term: is it a variable, what is its
name if it is one, and how is it written?

\begin{newfile}{terms.nex}
-- terms.nex: what Datalog talks about. Terms, atoms, literals, rules,
-- and whole programs.

-- A term is a constant (a symbol, an integer or a string) or a variable.
union Term
  Symbol(name: String)
  Int(value: Integer)
  Text(value: String)
  Var(name: String)
end

function is_variable(t: Term): Boolean do
  match t of
    Var then result := true
    _   then result := false
  end
end

-- A variable's name, or "" for a constant.
function variable_name(t: Term): String do
  match t of
    Var(name) then result := name
    _         then result := ""
  end
end

-- A term as Datalog writes it.
function term_text(t: Term): String do
  match t of
    Symbol(name) then result := name
    Int(value)   then result := value.to_string
    Text(value)  then result := "\"" + value + "\""
    Var(name)    then result := name
  end
end
\end{newfile}

\code{variable_name} returns \code{""} for a constant, so that code that
only cares about variables can skip constants with one test instead of a
\code{match}.

\section{Atoms}

An atom is a predicate and a list of terms. Its \code{text} is the atom as
Datalog writes it, and it does double duty: two atoms are the same exactly
when their texts are, so the text can serve as a key when we store facts
in Chapter~\ref{ch:database}. That works because \code{term_text} writes
every kind of constant differently: a string has quotes, a symbol does not,
so \code{alice} and \code{"alice"} never collide.

\begin{continued}{terms.nex}
-- A predicate applied to terms: parent(alice, bob), ancestor(X, Y).
class Atom
create
  make(p: String, a: Array[Term]) do
    predicate := p
    args := a
  end
feature
  predicate: String
  args: Array[Term]

  arity: Integer do
    result := args.length
  end

  -- Does it mention no variables? Facts are ground atoms.
  is_ground: Boolean do
    result := true
    across args as t do
      if is_variable(t) then
        result := false
      end
    end
  end

  -- The names of its variables, each once, in the order they appear.
  variables: Array[String] do
    result := []
    across args as t do
      let name := variable_name(t)
      if name /= "" and not result.contains(name) then
        result.add(name)
      end
    end
  end

  -- The atom as Datalog writes it. Two atoms are the same exactly when
  -- their texts are, so the text also serves as a key.
  text: String do
    result := predicate
    if arity > 0 then
      result := result + "("
      from let i := 0 until i = arity do
        if i > 0 then
          result := result + ", "
        end
        result := result + term_text(args.get(i))
        i := i + 1
      end
      result := result + ")"
    end
  end
invariant
  named: predicate.length > 0
end
\end{continued}

\section{Literals and rules}

A literal is again one of a few kinds. We define all three now, although
the tests come into their own in Chapter~\ref{ch:joins} and negation in
Chapter~\ref{ch:negation}.

\begin{continued}{terms.nex}
-- One condition in the body of a rule.
union Literal
  Positive(atom: Atom)
  Compare(op: String, left: Term, right: Term)
  Negated(atom: Atom)
end

function literal_text(l: Literal): String do
  match l of
    Positive(atom) then result := atom.text
    Compare(op, left, right) then
      result := term_text(left) + " " + op + " " + term_text(right)
    Negated(atom) then result := "not " + atom.text
  end
end
\end{continued}

A rule has a head and a body. Its constructor requires a non-empty body:
a ``rule'' with nothing to check is a fact, and facts are stored
separately.

\begin{continued}{terms.nex}
-- head :- body. A rule with an empty body would be a fact.
class Rule
create
  make(h: Atom, b: Array[Literal])
  require
    has_body: b.length > 0
  do
    head := h
    body := b
  end
feature
  head: Atom
  body: Array[Literal]

  text: String do
    result := head.text + " :-"
    from let i := 0 until i = body.length do
      result := result + when i = 0 then " " else ", " end + literal_text(body.get(i))
      i := i + 1
    end
    result := result + "."
  end
end
\end{continued}

Finally, a whole program is its facts, its rules and its questions. A fact
must be ground, and \code{add_fact} says so: a program that states
\code{parent(X, bob)} has made a mistake, and the parser will report it
before it gets this far.

\begin{continued}{terms.nex}
-- A whole program: the facts it states, the rules it gives, and the
-- questions it asks.
class Datalog_Program
create
  make do
    facts := []
    rules := []
    queries := []
  end
feature
  facts: Array[Atom]
  rules: Array[Rule]
  queries: Array[Atom]

  add_fact(a: Atom)
  require
    ground: a.is_ground
  do
    facts.add(a)
  ensure
    one_more: facts.length = old facts.length + 1
  end

  add_rule(r: Rule) do
    rules.add(r)
  end

  add_query(q: Atom) do
    queries.add(q)
  end
end
\end{continued}

\section{Trying it out}

There is no parser yet, so the try program builds a rule by hand: the
recursive rule for \code{ancestor}. Each term is declared as a \code{Term};
without that, \code{[x, z]} would be an array of \code{Var}, the most
specific type its elements share, and \code{Atom.make} wants an array of
\code{Term}.

\begin{tryit}{try_terms.nex}
intern Terms

-- Each term is declared a Term, so that arrays of them are Array[Term].
let x: Term := create Var.make("X")
let y: Term := create Var.make("Y")
let z: Term := create Var.make("Z")
let head := create Atom.make("ancestor", [x, z])
let parent := create Atom.make("parent", [x, y])
let rest := create Atom.make("ancestor", [y, z])
let first: Literal := create Positive.make(parent)
let second: Literal := create Positive.make(rest)
let rule := create Rule.make(head, [first, second])
console.print_line(rule.text)
console.print_line(parent.variables.to_string)
console.print_line(parent.is_ground.to_string)

let alice: Term := create Symbol.make("alice")
let ann: Term := create Text.make("Ann Lee")
let fact := create Atom.make("parent", [alice, ann])
console.print_line(fact.text)
console.print_line(fact.is_ground.to_string)
\end{tryit}
\tryoutput{try_terms}
